%==============================================================
%  COMP4328/5328/8328 Advanced Machine Learning -- Assignment 1
%
%  Layout follows the brief exactly:
%     Font  : Times New Roman  (mathptmx / newtxtext give the Times face
%                               that TeX distributions ship as the metric-
%                               compatible clone of Times New Roman)
%     Title : 14 pt
%     Body  : 12 pt
%     Length: 10-15 pages, hard maximum 20
%
%  Build:   cd report && latexmk -pdf report.tex      (or make)
%  Figures are read straight out of ../code/results/, so re-running
%  code/run_all.sh and rebuilding is enough to refresh every plot.
%==============================================================
\documentclass[12pt,a4paper]{article}

% ---------- fonts: Times New Roman ----------
\usepackage[T1]{fontenc}
\usepackage[utf8]{inputenc}
\usepackage{mathptmx}          % Times text + Times math
% If your TeX installation has the newer newtx bundle (MacTeX, TeX Live full,
% Overleaf all do), these two lines give a slightly better Times maths face.
% Comment out \usepackage{mathptmx} above before enabling them.
%   \usepackage{newtxtext}
%   \usepackage{newtxmath}

% ---------- layout ----------
\usepackage[a4paper,margin=0.85in]{geometry}
\usepackage{setspace}
\usepackage{parskip}

% Section headings kept close to body size so the document reads like a
% paper rather than a poster.  titlesec is the clean way to do this, but a
% minimal TeX install (BasicTeX, tinytex) may not ship it, so we fall back
% to plain LaTeX's own sectioning hooks instead of failing to build.
\IfFileExists{titlesec.sty}{%
  \usepackage{titlesec}%
  \titleformat{\section}{\normalfont\fontsize{13}{15}\bfseries}{\thesection}{1em}{}%
  \titleformat{\subsection}{\normalfont\fontsize{12}{14}\bfseries}{\thesubsection}{1em}{}%
}{%
  \makeatletter
  \renewcommand\section{\@startsection{section}{1}{\z@}%
    {-3.0ex \@plus -1ex \@minus -.2ex}{2.0ex \@plus.2ex}%
    {\normalfont\fontsize{13}{15}\bfseries}}
  \renewcommand\subsection{\@startsection{subsection}{2}{\z@}%
    {-2.5ex \@plus -1ex \@minus -.2ex}{1.5ex \@plus.2ex}%
    {\normalfont\fontsize{12}{14}\bfseries}}
  \makeatother
}

% ---------- maths, figures, tables ----------
\usepackage{amsmath,amssymb,bm}
\usepackage{graphicx}
\usepackage{booktabs}
\IfFileExists{multirow.sty}{\usepackage{multirow}}{}
\usepackage{subcaption}
\usepackage{float}
\usepackage[table]{xcolor}
% algorithm/algpseudocode are only used for the pseudocode floats; a minimal
% TeX install may lack them, and the report still builds without them.
\IfFileExists{algorithm.sty}{\usepackage{algorithm}}{}
\IfFileExists{algpseudocode.sty}{\usepackage{algpseudocode}}{}
\usepackage[numbers,sort&compress]{natbib}
\usepackage[colorlinks=true,allcolors=blue]{hyperref}
\usepackage{url}

\graphicspath{{../code/results/}{figures/}}

% ---------- compactness: the brief caps the report at 20 pages ----------
% Font sizes stay exactly as required (14 pt title, 12 pt body); only the
% whitespace around floats and paragraphs is tightened.
\setlength{\parskip}{0.35\baselineskip}
\usepackage{caption}
\captionsetup{font=small,skip=4pt}
\setlength{\textfloatsep}{8pt plus 2pt minus 2pt}
\setlength{\intextsep}{8pt plus 2pt minus 2pt}
\setlength{\abovecaptionskip}{4pt}
\setlength{\belowcaptionskip}{0pt}
% float placement: allow floats to fill most of a page alongside text, so that
% LaTeX does not spend whole pages on floats alone.
\renewcommand{\topfraction}{0.92}
\renewcommand{\bottomfraction}{0.75}
\renewcommand{\textfraction}{0.06}
\renewcommand{\floatpagefraction}{0.85}
\setcounter{topnumber}{3}
\setcounter{bottomnumber}{2}
\setcounter{totalnumber}{5}

% ---------- drafting helper: delete before submitting ----------
\newcommand{\todo}[1]{\textcolor{red}{\textbf{#1}}}
% Turn the markers off for a clean build:  \renewcommand{\todo}[1]{}


% ---------- figure helper ----------
% Includes a figure from ../code/results/ if it exists, and otherwise draws a
% labelled placeholder box.  This means report.tex always compiles, even
% before the (multi-hour) Extended YaleB experiments have finished -- you can
% write and check the layout while the numbers are still being produced.
\newcommand{\resultfig}[2][\linewidth]{%
  \IfFileExists{../code/results/#2}%
    {\includegraphics[width=#1]{#2}}%
    {\fbox{\parbox[c][0.26\linewidth][c]{\dimexpr#1-2\fboxsep-2\fboxrule\relax}%
      {\centering\footnotesize\texttt{\detokenize{#2}}\\[4pt]
       \itshape not generated yet --- run \texttt{code/run\char`_all.sh}}}}%
}

% ---------- table helper ----------
% \input a table generated by code/make_tables.py, silently skipping it if the
% corresponding experiment has not been run yet.  Regenerate them all with
%     cd code && python make_tables.py
\newcommand{\resulttable}[1]{\IfFileExists{tables/#1.tex}{\input{tables/#1.tex}}{}}

% ---------- shorthands ----------
\newcommand{\bV}{\bm{V}}
\newcommand{\bW}{\bm{W}}
\newcommand{\bH}{\bm{H}}
\newcommand{\bS}{\bm{S}}
\newcommand{\bD}{\bm{D}}
\newcommand{\bR}{\bm{R}}
\newcommand{\bT}{\bm{T}}
\newcommand{\bK}{\bm{K}}
\newcommand{\bL}{\bm{L}}
\newcommand{\bv}{\bm{v}}
\newcommand{\bw}{\bm{w}}
\newcommand{\br}{\bm{r}}
\newcommand{\Frob}[1]{\left\lVert #1 \right\rVert_F}
\newcommand{\Lone}[1]{\left\lVert #1 \right\rVert_1}
\newcommand{\Ltwoone}[1]{\left\lVert #1 \right\rVert_{2,1}}
\newcommand{\RRE}{\mathrm{RRE}}

%==============================================================
\begin{document}

% ---------------- title block (14 pt) ----------------
\begin{center}
  {\fontsize{14}{17}\selectfont\bfseries
   Robustness of Non-negative Matrix Factorization\\
   under Contiguous Occlusion Noise\par}
  \vspace{0.6em}
  {\fontsize{12}{14}\selectfont
   COMP4328/5328/8328 Advanced Machine Learning --- Assignment 1\par}
  \vspace{0.8em}
  {\fontsize{12}{14}\selectfont
   \textbf{Group ID:} \todo{XX}\par}
  \vspace{0.4em}
  {\fontsize{11}{13}\selectfont
   \begin{tabular}{lll}
     \todo{Name 1} & \todo{SID 1} & \todo{unikey 1} \\
     \todo{Name 2} & \todo{SID 2} & \todo{unikey 2} \\
     \todo{Name 3} & \todo{SID 3} & \todo{unikey 3} \\
     \todo{Name 4} & \todo{SID 4} & \todo{unikey 4} \\
   \end{tabular}\par}
\end{center}
\vspace{0.6em}

%==============================================================
\begin{abstract}
\noindent
% Rubric: "problem, methods, organization" [3 marks].
% Keep to ONE paragraph. Write this LAST, once the numbers are final.
% A working template with the four moves an abstract needs:
%   (1) problem  (2) what we did  (3) headline result  (4) organisation.
Non-negative Matrix Factorization (NMF) learns a parts-based, additive
representation of non-negative data, but its standard Frobenius-norm
formulation is a least-squares estimator and therefore has an unbounded
influence function: a handful of grossly corrupted pixels can dominate the
objective and destroy the learned basis. This report studies how far that
weakness can be repaired by changing the loss. We implement five NMF
algorithms from scratch in NumPy --- the Frobenius-norm multiplicative-update
baseline, $L_1$-norm NMF, hypersurface-cost NMF, $L_{2,1}$-norm NMF, and an
$L_1$-regularised robust NMF with an explicit sparse outlier term --- and
evaluate them on the ORL and Extended YaleB face datasets contaminated by
contiguous occlusion blocks. We vary three independent properties of the
corruption: the block edge length, the number of blocks per image, and the
\emph{proportion of images} corrupted, and measure relative reconstruction
error, clustering accuracy and normalised mutual information over five
independent $90\%$ subsamples. The sparse-outlier model cuts
RRE by $52\%$ relative to the Frobenius baseline at $b=15$ on Extended YaleB
($0.323$ against $0.668$), while $L_{2,1}$-NMF is numerically indistinguishable
from that same baseline --- agreeing to three decimal places --- on every
condition in which all images are corrupted equally, and separates from it only
once the corruption is unevenly distributed across samples. Section~\ref{sec:intro} introduces the problem,
Section~\ref{sec:related} reviews related NMF variants,
Section~\ref{sec:methods} derives the five algorithms and the noise model,
Section~\ref{sec:experiments} reports the experiments, and
Section~\ref{sec:conclusion} concludes.
\end{abstract}

%==============================================================
\section{Introduction}
\label{sec:intro}
% Rubric [5 marks]: "What is the problem you intend to solve? Why is this
% problem important?" -- plus the brief asks for the main idea of NMF, its
% applications, and an overview of the methods used.

\subsection{Non-negative matrix factorization}
% WHAT TO WRITE:
%  - the factorisation V ~ WH, shapes, and the meaning of the two factors
%  - why non-negativity gives PARTS-BASED representations, contrast with PCA
%    (which allows cancellation between positive and negative components)
%  - applications: face/object representation, document topic models,
%    hyperspectral unmixing, audio source separation, recommender systems
Given a non-negative data matrix $\bV \in \mathbb{R}^{m \times n}_{+}$ whose
$n$ columns are vectorised images of $m$ pixels each, NMF seeks factors
$\bW \in \mathbb{R}^{m \times k}_{+}$ and $\bH \in \mathbb{R}^{k \times n}_{+}$
with $k \ll \min(m,n)$ such that $\bV \approx \bW\bH$. The two factors admit
a direct reading: the $k$ columns of $\bW$ are themselves images and act as a
learned dictionary of \emph{basis images}, while column $j$ of $\bH$ holds the
$k$ coefficients that reconstruct image $j$ as the additive mixture
$\bv_j \approx \sum_{t=1}^{k} \bw_t H_{tj}$.

What distinguishes NMF from the other classical low-rank factorisations is
precisely the constraint that makes it harder to solve. Principal component
analysis minimises the same reconstruction error over unconstrained factors,
so its components may take either sign and a reconstruction may \emph{cancel}:
a bright region in one component can be removed by a negative coefficient on
another. The individual components are consequently holistic --- eigenfaces
look like whole, ghostly faces --- and are not interpretable as parts. Forbidding
negative entries forbids cancellation outright. Every basis image can only ever
be \emph{added} to the reconstruction, so the only way to represent a set of
faces with few components is to discover pieces that actually recur across
them, which is why NMF bases for face data tend to localise onto eyes, noses,
mouths and hairlines \citep{lee1999learning}. The same argument explains the
breadth of its applications: topics as additive mixtures of words in document
collections, endmembers in hyperspectral unmixing, spectrogram atoms in audio
source separation, and latent factors in recommendation, all of which are
domains where the data are non-negative and a subtractive explanation would be
physically meaningless.

\subsection{The problem: robustness to gross corruption}
% WHAT TO WRITE:
%  - the standard objective is ||V - WH||_F^2, i.e. a Gaussian/MLE assumption
%  - real face data violates it: occlusions, shadows, sensor faults are
%    SPARSE but LARGE, not small and Gaussian
%  - a squared loss has an unbounded influence function -> breakdown point 0
%  - therefore robust losses / explicit outlier models are needed
%  - state the concrete question this report answers
The standard NMF objective is the squared Frobenius error
$\|\bV - \bW\bH\|_F^2$. Minimising it is maximum-likelihood estimation under
the assumption that every entry of $\bV$ is the true value plus independent
Gaussian noise of a common variance. That assumption is a statement about the
\emph{whole} data matrix, and real face imagery violates it in a specific and
damaging way. Sunglasses, scarves, a hand in front of a face, specular
highlights, cast shadows, compression blocks and dead sensor pixels do not
perturb many pixels slightly; they replace a \emph{few} pixels
\emph{completely}. The corruption is sparse in its support but unbounded in
magnitude --- the opposite of the Gaussian regime.

Under a squared loss the consequences are not merely a loss of accuracy but a
loss of the estimator itself. The influence function of the least-squares
loss, $\psi(r) = \partial_r r^2 = 2r$, is unbounded, so the contribution of a
single residual to the gradient grows without limit as that residual grows. An
occlusion block of edge length $b$ filled with the saturated value contributes
on the order of $b^2$ squared residuals of near-maximal size, and the optimiser
will happily sacrifice the fit on thousands of clean pixels to reduce them. In
the language of robust statistics the breakdown point is zero: an arbitrarily
small fraction of arbitrarily corrupted entries suffices to drive the estimate
arbitrarily far from the truth. Empirically this shows up as bright,
square-shaped artefacts absorbed directly into the learned basis images, which
is visible in Section~\ref{sec:res-qualitative}.

Two families of remedy exist. One may \emph{change the loss} so that the
influence of a large residual is bounded --- replacing the square by an
absolute value ($L_1$), by a smooth interpolation between the two
(hypersurface / pseudo-Huber), or by a per-column norm that down-weights whole
corrupted samples ($L_{2,1}$). Alternatively one may \emph{model the
corruption explicitly}, adding a sparse outlier matrix that absorbs the
occlusion so that the low-rank part never has to explain it. This report asks a
concrete question about these remedies: given contiguous occlusion noise of a
controlled size, count and prevalence, which of them actually recovers the
clean signal, by how much, and --- the question that turns out to matter most
--- \emph{under which corruption geometry} does each one help at all?

\subsection{Contributions}
This report makes the following contributions.
\begin{enumerate}\itemsep0pt
  \item We implement five NMF algorithms from scratch using only NumPy and
        the Python standard library, and derive all five as instances of a
        single majorise--minimise (MM) template, which isolates their only
        real difference: \emph{at what granularity each one down-weights a
        residual} (Section~\ref{sec:methods}).
  \item We identify and correct a subtlety that is easy to miss in this
        experimental design: a column-wise robust loss such as
        $L_{2,1}$ can only distinguish clean from corrupted samples if the
        corruption is \emph{unevenly distributed across samples}. We
        therefore add the contamination \emph{ratio} as a third noise axis
        (Section~\ref{sec:noise}).
  \item We report a rigorous evaluation --- three metrics, five independent
        $90\%$ subsamples, and an explicit convergence study --- on both
        required datasets (Section~\ref{sec:experiments}).
\end{enumerate}

%==============================================================
\section{Related work}
\label{sec:related}
% Rubric [6 marks]: "Previous relevant methods used in literature?"
% The brief asks for the MAIN IDEA of related NMF algorithms INCLUDING
% THEIR ADVANTAGES AND DISADVANTAGES. Marks are lost for a list of names
% with no critique, so give each paragraph an explicit "advantage /
% disadvantage" sentence.

\subsection{Classical NMF}
\citet{paatero1994positive} introduced positive matrix factorization, and
\citet{lee1999learning} popularised NMF together with the multiplicative
update (MU) rules that make the non-negativity constraint free of charge.
Their argument is an auxiliary-function (majorise--minimise) one: for the
current iterate one constructs a surrogate $G(\bH, \bH^{t})$ that lies above
the objective everywhere and touches it at $\bH^{t}$, and then minimises the
surrogate in closed form. Because $F(\bH^{t+1}) \le G(\bH^{t+1}, \bH^{t}) \le
G(\bH^{t}, \bH^{t}) = F(\bH^{t})$, the objective is non-increasing at every
step by construction, and because the update is a \emph{ratio} of non-negative
quantities the iterates stay feasible without any projection or step-size
search. \citet{lee2001algorithms} give the same treatment for the
Kullback--Leibler divergence.

\emph{Advantages.} No step size to tune, no constraint to enforce, $O(mnk)$ per
update, and monotone descent that makes the method trivial to debug.
\emph{Disadvantages.} Monotone descent is weaker than convergence:
\citet{lin2007convergence} showed the original rules need not reach a
stationary point, and the tail convergence is slow. More importantly here, both
classical costs are maximum-likelihood estimators for a \emph{specific} noise
model --- Gaussian for the Frobenius norm, Poisson for KL --- and both give
large residuals unbounded influence.

\subsection{Robust loss functions}
The first family keeps the factorisation model and replaces the discrepancy
measure.

\paragraph{$L_1$-norm NMF.}
\citet{ke2005robust} and \citet{ding2006orthogonal} replace the squared error by
the entrywise $L_1$ norm, the maximum-likelihood criterion under Laplacian
noise. Its influence function $\mathrm{sign}(r)$ is bounded, so a pixel off by
$0.9$ pulls no harder than one off by $0.1$ and the fit can abandon occluded
pixels instead of chasing them.
\emph{Advantage:} robustness at the granularity of the individual pixel, which
matches the support of occlusion noise exactly.
\emph{Disadvantage:} non-differentiable at zero, and the reweighted surrogate's
$1/|r|$ weights blow up on the well-fitted pixels, so a smoothing constant
$\delta$ must be introduced --- a genuine hyper-parameter
(Section~\ref{sec:sensitivity}).

\paragraph{Hypersurface cost.}
\citet{hamza2006reconstruction} propose the pseudo-Huber loss
$\phi(r) = \sqrt{1 + r^{2}} - 1$, quadratic for $|r| \ll 1$ and linear for
$|r| \gg 1$.
\emph{Advantage:} differentiable everywhere, so no smoothing constant is
needed, and it combines quadratic efficiency on inliers with $L_1$'s bounded
influence on outliers.
\emph{Disadvantage:} the crossover sits at $|r| = 1$ in the units of the data,
so with pixels in $[0,1]$ every residual stays in the quadratic regime and the
method silently degenerates into the Frobenius baseline unless a scale constant
$c$ is introduced --- a trap we examine in
Section~\ref{sec:alg-hypersurface}.

\paragraph{$L_{2,1}$-norm NMF.}
\citet{kong2011robust} use $\|\bR\|_{2,1} = \sum_{j} \|\br_j\|_2$, the sum over
columns of the per-column Euclidean norm. Each column is one image, so the
cost grows only \emph{linearly} in the size of an image's residual, and the
induced reweighting $d_j = 1/\|\br_j\|_2$ acts as a soft sample-selection
mechanism that discounts badly damaged images as a whole.
\emph{Advantage:} the right model when entire observations are unreliable ---
a corrupted photograph, a mislabelled document, a faulty sensor channel.
\emph{Disadvantage, and the one that drives the experimental design of this
report:} the robustness is only at the sample level. If \emph{every} column is
corrupted to the same degree, all $d_j$ are nearly equal, a uniform column
weighting is absorbed by the scaling freedom of the factorisation, and
$L_{2,1}$-NMF reduces numerically to plain Frobenius NMF. It cannot see
\emph{which pixels} inside a column are bad, only \emph{which columns} are
worse than the others.

\subsection{Explicit outlier models}
A second family leaves the loss quadratic and instead enlarges the model.
Robust PCA \citep{candes2011robust} decomposes $\bV = \bL + \bS$ into a
low-rank term and a sparse term and recovers both exactly, under incoherence
conditions, by minimising $\|\bL\|_* + \lambda\|\bS\|_1$.
\citet{zhang2011robust} carry the same idea into the non-negative setting,
writing $\bV \approx \bW\bH + \bS$ with an $L_1$ penalty on $\bS$ and solving
by alternating soft-thresholding of $\bS$ with multiplicative updates of
$\bW, \bH$; \citet{ding2010rpca} give a related non-negative treatment.

\emph{Advantages.} The corruption is \emph{explained} rather than discounted:
the low-rank part is never asked to represent the occlusion, which is the right
prior for a contiguous white block, and the recovered $\bS$ is an interpretable
occlusion map. The remaining loss is quadratic, so Lee--Seung applies unchanged
to $\bW$ and $\bH$.
\emph{Disadvantages.} An extra $m \times n$ variable doubles the memory
footprint, and the penalty $\lambda$ trades reconstruction against sparsity ---
too small and $\bS$ absorbs the signal, too large and it stays empty
(Section~\ref{sec:sensitivity}).

\subsection{Other directions}
A third line of work adds \emph{structure} to the representation rather than
robustness to the loss: graph-regularised NMF \citep{cai2011graph}, sparse NMF
\citep{hoyer2004nonnegative}, orthogonal NMF \citep{ding2006orthogonal} and
deep multi-layer NMF. These are complementary to the methods studied here ---
they change what a \emph{good} factorisation looks like, whereas we ask whether
the factorisation survives grossly corrupted inputs at all. We therefore hold
the representation-side choices fixed (rank $k$ equal to the class count, no
extra regularisation) so that the only variable across the five algorithms is
how each treats a large residual.

%==============================================================
\section{Methods}
\label{sec:methods}
% Rubric [25 marks]: pre-processing; NMF algorithm formulation; noise
% choice, description and illustration. The brief additionally asks you to
% "explain the robustness of each algorithm from a theoretical view" -- the
% influence-function paragraph below is where those marks are.

\subsection{Notation and pre-processing}
\label{sec:preprocessing}
Each image is converted to 8-bit greyscale, downsampled by an integer factor
($3$ for ORL, giving $37 \times 30$ pixels; $4$ for Extended YaleB, giving
$48 \times 42$), vectorised column-wise, and divided by $255$ so that every
entry lies in $[0,1]$. Downsampling is the only lossy step and is applied for
tractability: it reduces the per-iteration cost of the multiplicative updates,
which is $O(mnk)$. The resulting data matrices are
$\bV \in \mathbb{R}^{1110 \times 400}_{+}$ for ORL ($40$ subjects, $10$ images
each) and $\bV \in \mathbb{R}^{2016 \times 2414}_{+}$ for Extended YaleB
($38$ subjects; the per-folder \texttt{*Ambient.pgm} frames, which are
illumination references rather than face captures, are excluded by the
loader). The rank is set to the number of classes present in the subsample,
$k = 40$ and $k = 38$ respectively, so that the comparison is never confounded
by a difference in model capacity.

Deliberately, \emph{no} further pre-processing is applied. Mean-centring and
whitening --- both standard before PCA --- are inadmissible here: each would
introduce negative entries and so violate the constraint that defines the
model. We likewise do not contrast-normalise the Extended YaleB images, even
though their illumination varies by orders of magnitude, because doing so
would suppress exactly the kind of large-magnitude structured variation whose
effect on the factorisation this report sets out to measure. The only
operations applied to the data are therefore greyscale conversion,
downsampling and a global division by $255$, all of which are order-preserving
and non-negativity-preserving.

\subsection{A single majorise--minimise template}
All algorithms except the sparse-outlier model of
Section~\ref{sec:l1reg} minimise a separable cost of the residual
$\bR = \bV - \bW\bH$,
\begin{equation}
  E(\bW,\bH) \;=\; \sum_{ij} \phi\!\left(R_{ij}\right),
  \qquad \bW,\bH \ge 0 ,
  \label{eq:generic}
\end{equation}
for a function $\phi$ that is concave and increasing in $r^2$. For any such
$\phi$, concavity gives the tangent upper bound
\begin{equation}
  \phi(r) \;\le\; \phi(r^{t}) \;+\;
  \underbrace{\frac{\phi'(r^{t})}{2 r^{t}}}_{=\,D^{t}}
  \left(r^{2} - (r^{t})^{2}\right),
  \label{eq:majoriser}
\end{equation}
which is a \emph{weighted quadratic} that touches $\phi$ at the current
iterate. Minimising the resulting weighted Frobenius surrogate with
Lee--Seung preconditioning yields the multiplicative updates
\begin{equation}
  \bW \leftarrow \bW \odot
    \frac{\left(\bD \odot \bV\right)\bH^{\top}}
         {\left(\bD \odot \bW\bH\right)\bH^{\top}},
  \qquad
  \bH \leftarrow \bH \odot
    \frac{\bW^{\top}\left(\bD \odot \bV\right)}
         {\bW^{\top}\left(\bD \odot \bW\bH\right)} ,
  \label{eq:mu}
\end{equation}
where $\odot$ and the fraction bar are elementwise and $\bD$ is held fixed
within an iteration and recomputed from the current residual at the start of
the next one (an iteratively reweighted least squares, or IRLS, scheme).
Because both updates are multiplicative and all factors are non-negative,
$\bW,\bH \ge 0$ is maintained without any projection step.

The practical consequence is that \textbf{the five algorithms differ only in
$\bD$}, as summarised in Table~\ref{tab:algorithms}. This is what makes the
comparison in Section~\ref{sec:experiments} a controlled experiment rather
than a comparison of five different codebases.

\begin{table}[tbp]
  \centering
  \caption{The five algorithms as instances of the same MM template.
           $\bR = \bV - \bW\bH$ is the current residual.}
  \label{tab:algorithms}
  \small
  \begin{tabular}{llll}
    \toprule
    Algorithm & Objective & Weight $\bD$ & Down-weights at \\
    \midrule
    F-norm MU      & $\Frob{\bR}^{2}$             & $1$                                  & --- (baseline) \\
    $L_1$-NMF      & $\Lone{\bR}$                 & $1/(|R_{ij}| + \delta)$              & pixel \\
    Hypersurface   & $\sum_{ij}\!\big(\!\sqrt{1+R_{ij}^{2}}-1\big)$ & $1/\sqrt{1+R_{ij}^{2}}$ & pixel \\
    $L_{2,1}$-NMF  & $\Ltwoone{\bR}$              & $1/\lVert \bR_{:j}\rVert_{2}$        & image \\
    $L_1$-Reg Robust & $\tfrac12\Frob{\bV-\bW\bH-\bS}^{2} + \lambda\Lone{\bS}$ & \multicolumn{2}{l}{explicit outlier variable $\bS$} \\
    \bottomrule
  \end{tabular}
\end{table}

\subsection{Algorithm 1: Frobenius-norm NMF (baseline, taught in this course)}
\label{sec:alg-frobenius}
\begin{equation}
  \min_{\bW,\bH \ge 0} \; \Frob{\bV - \bW\bH}^{2},
  \qquad \phi(r) = r^{2}, \quad D_{ij} = 1 .
\end{equation}
With $\bD = \mathbf{1}$ the template of Eq.~\eqref{eq:mu} collapses to the
classical Lee--Seung multiplicative updates
\begin{equation}
  \bW \leftarrow \bW \odot \frac{\bV\bH^{\top}}{\bW\bH\bH^{\top}},
  \qquad
  \bH \leftarrow \bH \odot \frac{\bW^{\top}\bV}{\bW^{\top}\bW\bH} .
\end{equation}
\emph{Monotonicity.} For fixed $\bW$, the function
$G(\bH,\bH^{t}) = F(\bH^{t}) + (\bH-\bH^{t})^{\top}\nabla F(\bH^{t})
+ \tfrac{1}{2}(\bH-\bH^{t})^{\top}\bK(\bH^{t})(\bH-\bH^{t})$, with the
diagonal matrix $K_{ab}(\bH^{t}) = \delta_{ab}(\bW^{\top}\bW\bH^{t})_{a}/H^{t}_{a}$,
majorises $F$ and touches it at $\bH^{t}$, because
$\bK(\bH^{t}) - \bW^{\top}\bW \succeq 0$ \citep{lee2001algorithms}. Setting
$\nabla_{\bH} G = 0$ yields precisely the update above, so
$F(\bH^{t+1}) \le G(\bH^{t+1},\bH^{t}) \le G(\bH^{t},\bH^{t}) = F(\bH^{t})$:
the objective never increases, and since it is bounded below by zero the
sequence of objective values converges. The same argument applies to $\bW$ by
symmetry.

\emph{Theoretical robustness: none.} The influence of a residual is
$\phi'(r) = 2r$, which is \emph{unbounded}. A single saturated pixel therefore
exerts arbitrarily large leverage on the gradient, and the estimator has a
breakdown point of zero. Concretely, a block of edge length $b$ filled with
$1.0$ contributes $O(b^{2})$ residuals of near-maximal magnitude, and since
each is squared the optimiser will trade away the fit on many clean pixels to
reduce them --- which is why the baseline's learned basis absorbs the
occlusion outright (Section~\ref{sec:res-qualitative}). Every other algorithm
below is a way of making $\phi'$ bounded, or of removing the offending
residuals from the loss altogether.

\subsection{Algorithm 2: $L_1$-norm NMF (not taught)}
\begin{equation}
  \min_{\bW,\bH \ge 0} \; \Lone{\bV - \bW\bH}
  = \sum_{ij} \left| V_{ij} - (\bW\bH)_{ij} \right| .
\end{equation}
Majorising $|r| \le r^{2}/(2s) + s/2$ at $s = |r^{t}|$ gives
$D_{ij} = 1/|R_{ij}|$, i.e.\ an iteratively-reweighted least-squares scheme in
which each pixel is weighted by the reciprocal of its own current error.

\emph{Theoretical robustness.} The influence function is
$\phi'(r) = \mathrm{sign}(r)$, bounded by $1$ regardless of $|r|$. A pixel
that is wrong by $0.9$ pulls on the solution no harder than one wrong by
$0.1$, so the fit is free to abandon the occluded pixels rather than chase
them. Because the weighting is \emph{entrywise}, the robustness operates at
exactly the granularity of the corruption: occlusion is sparse in the pixel
domain, and $L_1$-NMF discounts pixels.

\emph{The smoothing constant is not cosmetic.} The weight $1/|R_{ij}|$
diverges as a residual approaches zero --- which is what happens on the
well-fitted majority of pixels --- so in practice one uses
$D_{ij} = 1/(|R_{ij}| + \delta)$. This is not a numerical guard bolted onto
the $L_1$ scheme; it is exactly the IRLS weighting of the \emph{smoothed} loss
$\sqrt{r^{2}+\delta^{2}}$, so $\delta$ sets the residual scale below which the
loss is treated as quadratic. With pixels in $[0,1]$, our default
$\delta = 10^{-3}$ declares residuals under one tenth of one percent of the
dynamic range to be noise-free. Because this is a modelling choice rather than
a free knob, Section~\ref{sec:sensitivity} sweeps $\delta$ over four orders
of magnitude to confirm that the conclusions do not rest on the particular
value.

\subsection{Algorithm 3: hypersurface-cost NMF (not taught)}
\label{sec:alg-hypersurface}
\begin{equation}
  \min_{\bW,\bH \ge 0} \; \sum_{ij}
  \left( \sqrt{1 + \left(V_{ij}-(\bW\bH)_{ij}\right)^{2}} - 1 \right),
  \qquad D_{ij} = \frac{1}{\sqrt{1+R_{ij}^{2}}} .
\end{equation}
A Taylor expansion shows what this loss does: $\phi(r) \approx r^{2}/2$ for
$|r| \ll 1$ and $\phi(r) \approx |r|$ for $|r| \gg 1$. It is thus quadratic
--- and therefore statistically efficient --- on the clean pixels, and linear
--- and therefore robust --- on the corrupted ones, with a smooth transition
between the two. Its influence function $\phi'(r) = r/\sqrt{1+r^{2}}$ is
bounded by $1$, as for $L_1$, but unlike $L_1$ it is continuous at the origin,
so the loss is differentiable everywhere and needs \emph{no} hand-chosen
smoothing constant to be well posed \citep{hamza2006reconstruction}. Its
weights lie in $(0,1]$ rather than being unbounded, which makes the
reweighting markedly less aggressive than the $L_1$ scheme.

\paragraph{The scale constant, and a trap.} The transition between the two
regimes sits at $|r| = 1$ in the units in which the cost is written. Our
pixels live in $[0,1]$, so \emph{every attainable residual} satisfies
$|r| \le 1$ and the method never leaves its quadratic regime: with the
literature's $c = 1$ the hypersurface cost silently degenerates into the
Frobenius baseline, and its "robustness" is an illusion. We therefore
implement the scaled cost
$\phi_c(r) = c^{2}\big(\sqrt{1 + (r/c)^{2}} - 1\big)$ with weights
$D_{ij} = (1 + (R_{ij}/c)^{2})^{-1/2}$ and default $c = 0.05$, placing the
elbow at $5\%$ of the dynamic range so that clean pixels remain in the
quadratic regime while occluded ones (with $r$ of order $1$) are firmly in the
linear one. Section~\ref{sec:sensitivity} sweeps $c$ and shows the
degeneration happening as $c \to 1$.

\subsection{Algorithm 4: $L_{2,1}$-norm NMF (not taught)}
\begin{equation}
  \min_{\bW,\bH \ge 0} \; \Ltwoone{\bV - \bW\bH}
  = \sum_{j=1}^{n} \sqrt{\sum_{i=1}^{m} \left(V_{ij}-(\bW\bH)_{ij}\right)^{2}},
  \qquad d_{j} = \frac{1}{\lVert \bR_{:j}\rVert_{2}} .
\end{equation}
Since $\lVert \bm{r}\rVert_{2} = \min_{d>0}\big(\tfrac{d}{2}\lVert
\bm{r}\rVert_{2}^{2} + \tfrac{1}{2d}\big)$ with optimum
$d = 1/\lVert\bm{r}\rVert_{2}$, the surrogate is a \emph{column}-weighted
Frobenius cost and Eq.~\eqref{eq:mu} applies with $\bD = \mathbf{1}\bm{d}^{\top}$.

\paragraph{A prediction we test explicitly.} Because $d_j$ varies only
\emph{across} columns, $L_{2,1}$-NMF performs soft sample selection: it can
discount an image, never a pixel. If every image receives the same amount of
occlusion, all $\lVert\bR_{:j}\rVert_2$ are nearly equal, $\bm{d}$ is nearly
constant, and a constant column weighting is absorbed by the scale freedom of
the factorisation --- so $L_{2,1}$-NMF should then be numerically
\emph{indistinguishable} from the Frobenius baseline. This prediction is the
reason we sweep the contamination ratio in Section~\ref{sec:noise}, and
Section~\ref{sec:res-contamination} reports the test.

\subsection{Algorithm 5: $L_1$-regularised robust NMF (not taught)}
\label{sec:l1reg}
\begin{equation}
  \min_{\bW,\bH \ge 0,\; \bS} \;
  \tfrac{1}{2}\Frob{\bV - \bW\bH - \bS}^{2} + \lambda \Lone{\bS} .
\end{equation}
Rather than down-weighting corrupted entries, this model gives the corruption
its own variable: $\bS$ is an explicit sparse outlier matrix, so $\bW\bH$ is
fitted to the \emph{cleaned} data $\bV - \bS$. Block-coordinate descent
alternates a closed-form soft-thresholding step and a Frobenius MU step:
\begin{align}
  \bS &\leftarrow \mathcal{S}_{\lambda}\!\left(\bV - \bW\bH\right),
       &&\mathcal{S}_{\lambda}(x) = \mathrm{sign}(x)\max(|x|-\lambda, 0),
       \label{eq:sstep}\\
  \bW &\leftarrow \bW \odot \frac{\bT\bH^{\top}}{\bW\bH\bH^{\top}},
  \quad
  \bH \leftarrow \bH \odot \frac{\bW^{\top}\bT}{\bW^{\top}\bW\bH},
       &&\bT = \max(\bV - \bS, 0).
       \label{eq:whstep}
\end{align}
\emph{Descent.} Eq.~\eqref{eq:sstep} is the exact minimiser of the joint
objective in $\bS$ for fixed $(\bW,\bH)$ --- the proximal operator of
$\lambda\lVert\cdot\rVert_1$ --- and Eq.~\eqref{eq:whstep} is a Lee--Seung step
on the quadratic term with $\bS$ fixed, which is non-increasing by the
auxiliary-function argument of Section~\ref{sec:alg-frobenius}. Each block
step therefore leaves the joint objective non-increasing, and since it is
bounded below by $0$ the sequence of objective values converges. (As with all
block-coordinate schemes on a non-convex problem, this guarantees convergence
of the objective, not of the iterates to a global minimum.)

\emph{Choosing $\lambda$.} We use the standard Robust-PCA scaling
$\lambda = 1/\sqrt{\max(m,n)}$ \citep{candes2011robust}, which is the value for
which exact recovery of a sparse-plus-low-rank decomposition holds under
incoherence. It has an intuitive reading as a threshold: a residual must exceed
$\lambda$ before it is declared an outlier and moved into $\bS$. The sweep in
Section~\ref{sec:sensitivity} varies $\lambda$ by two orders of magnitude
around this value, and shows the two failure modes at the ends of the range
--- $\bS$ absorbing genuine signal when $\lambda$ is too small, and $\bS$
staying empty (so that the method reduces to the Frobenius baseline) when it is
too large.

\emph{Why this model should fit our noise best.} A white occlusion block is
\emph{precisely} a sparse, large-magnitude, additive perturbation: few entries
affected, each affected entry changed by an amount of order the full dynamic
range. That is the generative assumption of $\bV = \bW\bH + \bS$ stated
exactly. The other four algorithms can only \emph{discount} corrupted entries,
leaving a biased residual in the loss; this one \emph{removes} them from the
quadratic term altogether, and returns the estimated occlusion map $\bS$ as a
by-product that can be inspected directly. We should therefore expect it to
lead on the heavy-occlusion conditions, and Section~\ref{sec:experiments}
confirms that it does.

\subsection{Noise model}
\label{sec:noise}
Following the brief we use \emph{contiguous occlusion}: squares of edge length
$b$ filled with the saturated value $255$ (i.e.\ $1.0$ after normalisation)
are pasted at uniformly random positions. We control three independent
properties of the corruption:
\begin{enumerate}\itemsep0pt
  \item the \textbf{block edge length} $b \in \{0,5,10,15,20\}$ pixels
        ($b=0$ is the uncorrupted control);
  \item the \textbf{number of blocks} per image $\in \{0,1,2,3,4\}$ at $b=10$;
  \item the \textbf{contamination ratio}, i.e.\ the fraction of images that
        are corrupted at all, $\in \{0.1, 0.25, 0.5, 0.75, 1.0\}$ at $b=10$
        with one block.
\end{enumerate}
The first two axes are named in the brief. The third is ours, and it is not
cosmetic: axes 1--2 change \emph{how many pixels} are corrupted while keeping
every image equally damaged, whereas axis 3 changes \emph{how many images} are
corrupted. A per-pixel robust loss ($L_1$, hypersurface) responds to the
former; a per-image robust loss ($L_{2,1}$) can only respond to the latter.
Sweeping only axes 1--2 would therefore make $L_{2,1}$-NMF look identical to
the non-robust baseline and would hide the very phenomenon the assignment asks
us to demonstrate.

Because blocks are placed independently they may overlap, so the nominal
budget $n_{\text{blocks}} \cdot b^{2}$ overstates the damage; we therefore log
the \emph{measured} fraction of modified entries for every condition and
report it alongside the nominal settings.

\begin{figure}[tbp]
  \centering
  \begin{subfigure}{0.48\linewidth}
    \resultfig[\linewidth]{orl_noise_demo.png}
    \caption{ORL}
    \label{fig:orl-noise}
  \end{subfigure}
  \hfill
  \begin{subfigure}{0.48\linewidth}
    \resultfig[\linewidth]{yaleb_noise_demo.png}
    \caption{Extended YaleB}
    \label{fig:yaleb-noise}
  \end{subfigure}
  \caption{Demonstration of the noise: an original image and the same image
           under the occlusion conditions used in the experiments.}
  \label{fig:noise-demo}
\end{figure}

\subsection{Evaluation metrics}
\label{sec:metrics}
Let $\hat{\bV}$ be the clean data and $\bW,\bH$ the factors fitted on the
\emph{contaminated} matrix $\bV$. The required metric is the relative
reconstruction error
\begin{equation}
  \RRE = \frac{\Frob{\hat{\bV} - \bW\bH}}{\Frob{\hat{\bV}}} .
\end{equation}
Note the asymmetry: the algorithm sees only $\bV$ but is scored against
$\hat{\bV}$, so fitting the occlusion \emph{well} makes the score
\emph{worse}. We additionally report the two optional metrics, clustering
accuracy and normalised mutual information, obtained by running $K$-means with
$K$ equal to the number of classes on the columns of $\bH$ and mapping
clusters to labels by majority vote, as in the provided tutorial notebook. We
also log an accuracy computed with optimal one-to-one (Hungarian) matching,
which cannot be inflated by several clusters collapsing onto the same class.

\paragraph{Protocol.} Every experiment draws an independent uniform $90\%$
subsample of the images, and is repeated $5$ times; we report the mean and
the sample standard deviation across those repeats. Within a repeat, all noise
conditions and all algorithms see the identical subsample, and all algorithms
see the identical corrupted matrix, so differences between algorithms are not
confounded by the draw. The rank is fixed to $k = $ the number of classes
present in the subsample.

%==============================================================
\section{Experiments and discussion}
\label{sec:experiments}
% Rubric [25 marks]: experimental setup; results WITH COMMENTS; extensive
% analysis and discussion; relevant personal reflection. Marks here are for
% INTERPRETATION, not for more plots. For every figure, state (a) what the
% expected behaviour was, (b) what actually happened, (c) why.

\subsection{Experimental setup}
Every experiment factorises a contaminated matrix and scores the result
against the \emph{clean} one. The two datasets are summarised in
Table~\ref{tab:datasets}; all five algorithms of Section~\ref{sec:methods} are
run on every condition, so each cell of every table below is a like-for-like
comparison in which the only thing that differs between rows is the loss.

\paragraph{Noise grid.} Three axes are swept independently, each around a
common reference point. \emph{Block size:} $b \in \{0,5,10,15,20\}$ pixels
with one block per image, where $b=0$ is the uncontaminated control.
\emph{Block count:} $\{0,1,2,3,4\}$ blocks per image at a fixed $b=10$.
\emph{Contamination ratio:} the fraction of images corrupted,
$\{0.1, 0.25, 0.5, 0.75, 1.0\}$, at a fixed $b=15$ with $2$ blocks per
corrupted image. Blocks are square, filled with the saturated value $1.0$,
placed uniformly at random, and allowed to overlap.

\paragraph{Protocol.} Following the note in the brief, every condition is
repeated over $R=5$ independent random subsamples of $90\%$ of the images, and
we report mean and standard deviation across those repeats. The subsample for repeat $r$ is drawn once and \emph{shared by all
conditions and all algorithms} in that repeat, so differences between
algorithms are never confounded by differences in the data they saw. The rank
is $k = $ the number of classes present in the subsample ($40$ for ORL, $38$
for YaleB). Multiplicative updates run for at most $400$ iterations on ORL and
$250$ on YaleB, with an early stop when the relative change in the objective
falls below $\tau = 10^{-5}$; the convergence study in
Section~\ref{sec:convergence} justifies these budgets. Factors are initialised
from a fixed-seed uniform distribution, and the seed is a deterministic
function of (repeat, condition), so the entire grid is reproducible with
\texttt{bash run\_all.sh}.

\paragraph{Metrics.} All three metrics of Section~\ref{sec:metrics} --- RRE
(required), clustering accuracy and NMI (both optional) --- are reported for
every condition. Timings indicate relative cost between algorithms rather than
absolute performance.

\begin{table}[tbp]
  \centering
  \caption{Datasets after pre-processing.}
  \label{tab:datasets}
  \small
  \begin{tabular}{lccccc}
    \toprule
    Dataset & Images & Subjects & Native size & Reduce & $\bV$ \\
    \midrule
    ORL            & 400  & 40 & $112\times 92$  & 3 & $1110 \times 400$ \\
    Extended YaleB & 2414 & 38 & $192\times 168$ & 4 & $2016 \times 2414$ \\
    \bottomrule
  \end{tabular}
\end{table}

\subsection{Sanity check: the uncontaminated control}
Before any robustness claim can be made, the five algorithms must be shown to
agree when there is nothing to be robust to; a large gap at $b=0$ would
indicate an implementation or convergence fault rather than a property of the
loss. On clean ORL the Frobenius baseline attains
$\mathrm{RRE} = 0.1290 \pm 0.0005$ and $L_{2,1}$-NMF $0.1295 \pm 0.0007$ ---
indistinguishable, as Section~\ref{sec:methods} predicts they must be when no
column is worse than any other. Hypersurface-cost NMF is close behind at
$0.1388$, while $L_1$-NMF ($0.1849$) and the sparse-outlier model ($0.1711$)
are $43\%$ and $33\%$ worse respectively. The clustering metrics order the
five the same way (accuracy $0.760$ for the baseline against $0.641$ for
$L_1$), which confirms the ordering is a property of the fits and not an
artefact of the reconstruction metric.

That the robust losses are \emph{worse} on clean data is the usual
efficiency/robustness trade-off. Each achieves bounded influence by discounting
large residuals, and on clean data a large residual is not corruption but
signal --- a hard-to-reconstruct region of a face. Hypersurface-cost NMF pays
the smallest price because its weights are bounded in $(0,1]$, so its
reweighting is the gentlest. This row is the correct baseline for every later
table: a robust method must first recover this deficit before showing any net
benefit.

Extended YaleB reproduces the pattern with the gaps widened. The baseline
attains $0.1917 \pm 0.0009$ and $L_{2,1}$-NMF $0.1932 \pm 0.0005$ --- again
within a thousandth of each other --- while $L_1$-NMF ($0.2668$) and the
sparse-outlier model ($0.2935$) are $39\%$ and $53\%$ worse. The larger deficit
is expected here: YaleB's extreme illumination variation produces genuinely
large residuals in shadowed regions, and a robust loss cannot distinguish
"large because corrupted" from "large because dark". Clustering scores on YaleB
are low throughout (best accuracy $0.240$, against $0.760$ on ORL) for the same
reason --- a $K$-means partition of YaleB codes groups by illumination rather
than identity.

\subsection{Effect of occlusion size}
\label{sec:res-blocksize}
\begin{figure}[tbp]
  \centering
  \begin{subfigure}{0.46\linewidth}
    \resultfig[\linewidth]{orl_blocksize_rre.png}
    \caption{ORL}
  \end{subfigure}
  \hfill
  \begin{subfigure}{0.46\linewidth}
    \resultfig[\linewidth]{yaleb_blocksize_rre.png}
    \caption{Extended YaleB}
  \end{subfigure}
  \caption{RRE against occlusion block edge length $b$ (one block per image,
           all images corrupted). Error bars are one standard deviation over
           five $90\%$ subsamples.}
  \label{fig:blocksize}
\end{figure}
\resulttable{orl_blocksize_rre}
\resulttable{yaleb_blocksize_rre}
RRE grows monotonically with $b$ for all five algorithms, but not at the same
rate, and the ordering changes twice as $b$ increases.

At $b=5$ ($2.3\%$ of ORL pixels) only the hypersurface cost improves on the
baseline, $0.1443$ against $0.1846$, a $22\%$ reduction. $L_1$-NMF sits at
$0.1874$, no better than the baseline --- but it has degraded by $1\%$ from its
own $b=0$ value while the baseline degraded by $43\%$. It is not yet
\emph{ahead}; it has simply not moved.

At $b=10$ ($9.0\%$ of pixels) that difference in slope becomes decisive.
$L_1$-NMF reaches $0.2328$ and the sparse-outlier model $0.2358$ against
$0.3277$ --- a $29\%$ reduction --- and the clustering gap is larger still,
accuracy $0.491$ and $0.483$ against $0.271$. This is the regime the assignment
is about: corruption heavy enough to matter, sparse enough to be identifiable.

At $b=15$ and $b=20$ ($20\%$ and $36\%$ of pixels) the five converge again ---
$0.6154$ against $0.6363$ at $b=20$, a $3\%$ edge. Robustness is a statement
about a \emph{minority} of the data; once a third of every image is saturated
white the occlusion is a dominant, dataset-wide structure that a rank-$40$
model can simply represent. Every robust loss has a working range, and this
experiment locates it.

Throughout, $L_{2,1}$-NMF tracks the Frobenius baseline to three decimal places
at every single value of $b$ ($0.4696$ vs $0.4699$ at $b=15$;
$0.6364$ vs $0.6363$ at $b=20$). That is not a coding error --- it is exactly
the degeneracy predicted in Section~\ref{sec:methods}, because on this axis
\emph{every} image is corrupted identically and the column weights $d_j$ are
therefore near-uniform. Section~\ref{sec:res-contamination} tests that
explanation directly --- and the agreement is even tighter on YaleB, where
$L_{2,1}$ and the baseline differ by $0.0010$ at $b=15$ ($0.6665$ vs $0.6675$)
and by $0.0007$ at $b=20$ ($0.8955$ vs $0.8962$).

\paragraph{Extended YaleB tells the same story, but does not close at the top
end.} The advantage of the pixel-level robust losses is larger here and, unlike
on ORL, it keeps growing with $b$. At $b=10$, $L_1$-NMF reaches $0.2682$
against the baseline's $0.4502$, a $40\%$ reduction --- and that is within
$0.002$ of its own clean-data value ($0.2668$): raising the occlusion from
nothing to $5\%$ of all pixels cost $L_1$-NMF essentially nothing while it more
than doubled the baseline's error. At $b=15$ the sparse-outlier model takes the
lead with $0.3227$ against $0.6675$, a $52\%$ reduction, and at $b=20$ it still
delivers $0.5396$ against $0.8962$, a $40\%$ reduction, where on ORL the
corresponding gap had collapsed to $3\%$.

The contrast between datasets follows from the corrupted-pixel fractions, not
from the algorithms. YaleB images are larger ($48 \times 42$ against
$37 \times 30$), so the same $b=20$ destroys $36\%$ of an ORL image but only
$20\%$ of a YaleB one; the regime in which robustness stops helping is reached
on ORL and not on YaleB within the same nominal grid. At matched
corrupted-pixel fraction the two agree, which is what
Section~\ref{sec:res-blockcount} shows.

\subsection{Effect of the number of occlusions}
\label{sec:res-blockcount}
\begin{figure}[tbp]
  \centering
  \begin{subfigure}{0.46\linewidth}
    \resultfig[\linewidth]{orl_blockcount_rre.png}
    \caption{ORL}
  \end{subfigure}
  \hfill
  \begin{subfigure}{0.46\linewidth}
    \resultfig[\linewidth]{yaleb_blockcount_rre.png}
    \caption{Extended YaleB}
  \end{subfigure}
  \caption{RRE against the number of occlusion blocks per image ($b=10$).}
  \label{fig:blockcount}
\end{figure}
Adding blocks is qualitatively like enlarging them, and the ranking of the
five algorithms is the same as in Section~\ref{sec:res-blocksize}: the
sparse-outlier model leads at every non-zero count ($0.3552$ against $0.4331$
for the baseline at two blocks, a $18\%$ reduction), $L_1$-NMF is close behind,
hypersurface is a small improvement, and $L_{2,1}$ is again
indistinguishable from the baseline.

The more informative question is whether the two axes are the \emph{same}
experiment in disguise. Because we log the realised corrupted-pixel fraction
for every condition, the two axes can be compared at matched corruption
\emph{mass} rather than at matched nominal parameters. Table~\ref{tab:matched}
does this for ORL.

\begin{table}[tbp]
  \centering
  \caption{Two occlusion geometries compared at a matched corrupted-pixel
           fraction (ORL, mean RRE over five subsamples). Fewer large blocks
           and more small blocks do comparable damage.}
  \label{tab:matched}
  \small
  \begin{tabular}{llccc}
    \toprule
    Geometry & Corrupted pixels & F-norm MU & $L_1$-NMF & $L_1$-Reg Robust \\
    \midrule
    $b=10$, $4$ blocks & $29.4\%$ & $0.5682$ & $0.5448$ & $0.5141$ \\
    $b=15$, $2$ blocks & $33.7\%$ & $0.6076$ & $0.5905$ & $0.5775$ \\
    $b=20$, $1$ block  & $36.0\%$ & $0.6363$ & $0.6200$ & $0.6154$ \\
    \midrule
    $b=10$, $2$ blocks & $16.8\%$ & $0.4331$ & $0.3768$ & $0.3552$ \\
    $b=15$, $1$ block  & $20.3\%$ & $0.4699$ & $0.4530$ & $0.4466$ \\
    \bottomrule
  \end{tabular}
\end{table}

Read along a row-group, RRE is essentially a monotone function of the
corrupted-pixel fraction alone: the increase from $0.4331$ to $0.4699$ is
explained by the extra $3.5$ percentage points of corruption, not by the change
of geometry. The total \emph{mass} of corruption is what matters, not how it is
arranged. This is not trivial: one might expect a single large block, which
destroys a whole contiguous facial region, to be more damaging than the same
pixels scattered over four smaller ones. At this rank and these corruption
levels that effect is second order. It also simplifies the rest of the
study: the contamination-ratio axis of
Section~\ref{sec:res-contamination} can be read as varying \emph{who} is
corrupted while holding corruption mass per corrupted image fixed, and is
therefore genuinely a third, independent axis rather than a reparameterisation
of the first two.

Extended YaleB provides the cleanest possible version of this check. Two
geometries there land almost exactly on the same corrupted-pixel fraction by
different routes --- one block of edge $20$ destroys $19.8\%$ of the pixels,
two blocks of edge $15$ destroy $20.4\%$ --- and the baseline scores $0.8962$
and $0.8940$ on them respectively, a difference of $0.002$, well inside the
$\pm 0.004$ standard deviation. One big square and two medium ones are, to the
precision we can measure, the same experiment. The ranking is also unchanged
along the block-count axis on YaleB: the sparse-outlier model leads at every
non-zero count, from $0.3114$ against $0.6027$ at two blocks (a $48\%$
reduction) to $0.3473$ against $0.8072$ at four (a $57\%$ reduction), and
$L_{2,1}$ again shadows the baseline to within $0.002$ throughout.

\subsection{Effect of the contamination ratio}
\label{sec:res-contamination}
\begin{figure}[tbp]
  \centering
  \begin{subfigure}{0.46\linewidth}
    \resultfig[\linewidth]{orl_contamination_rre.png}
    \caption{ORL}
  \end{subfigure}
  \hfill
  \begin{subfigure}{0.46\linewidth}
    \resultfig[\linewidth]{yaleb_contamination_rre.png}
    \caption{Extended YaleB}
  \end{subfigure}
  \caption{RRE against the fraction of corrupted images. Corrupted images
           receive two blocks of edge length $b=15$; the remaining images are
           left clean.}
  \label{fig:contamination}
\end{figure}
\resulttable{orl_contamination_rre}
\resulttable{yaleb_contamination_rre}
This axis tests the prediction of Section~\ref{sec:methods} directly, and the
prediction holds in both directions.

At a contamination ratio of $0.1$ --- $36$ of $360$ ORL images corrupted, only
$3.2\%$ of all pixels --- $L_{2,1}$-NMF is the \emph{best} of the five,
$0.2209 \pm 0.0054$ against $0.2262 \pm 0.0055$ for the Frobenius baseline,
and it also attains the highest clustering accuracy ($0.626$ against $0.606$).
This is the only part of the entire study in which $L_{2,1}$ separates from the
baseline at all, and it is precisely the regime its column-wise reweighting was
designed for: $324$ columns are clean, $36$ are badly damaged, the weights
$d_j = 1/\|\br_j\|_2$ can tell them apart, and the damaged columns are
effectively discounted.

As the ratio rises the gap closes exactly as predicted. At $0.25$ the two are
$0.3188$ and $0.3235$; at $0.5$, $0.4386$ and $0.4418$; at $0.75$, $0.5311$ and
$0.5326$; and at $1.0$ they agree to three decimal places, $0.6074$ against
$0.6076$. The mechanism is not subtle: when every column is corrupted to the
same degree the weights $d_j$ become uniform, a uniform column scaling is
absorbed by the scale freedom of the factorisation, and $L_{2,1}$-NMF
\emph{is} Frobenius NMF up to numerical noise.

The methodological moral is the most transferable result in this report. The
protocol the brief literally describes --- corrupt every image and sweep $b$
--- is the one setting in which a column-wise robust loss provably cannot show
its advantage, and an experimenter using only that protocol would conclude,
wrongly, that $L_{2,1}$-NMF does not work. The two families also mirror each
other along this axis: $L_1$-NMF and the sparse-outlier model, which discount
\emph{pixels}, are worse than the baseline at ratio $0.1$ ($0.2462$ and
$0.2377$, their clean-data deficit still dominating at $3.2\%$ corruption) and
best at ratio $1.0$ ($0.5905$ and $0.5775$ against $0.6076$).

On Extended YaleB the $L_{2,1}$ signature is present but reads more clearly in
the clustering metrics than in RRE. At ratio $0.1$, $L_{2,1}$-NMF beats the
baseline on RRE ($0.3221$ against $0.3326$) and, more tellingly, attains the
best clustering of any algorithm on that condition --- accuracy $0.181$ and NMI
$0.283$, against $0.163$ and $0.260$ for the baseline. At ratio $1.0$ the two
coincide again, $0.8929$ against $0.8940$. The prediction therefore holds on
both datasets: $L_{2,1}$ separates from the baseline exactly where the
corruption is unevenly distributed, and nowhere else.

What differs on YaleB is that $L_{2,1}$ does not \emph{win} that condition
outright: the sparse-outlier model is already ahead on RRE ($0.3007$) and
extends the lead as the ratio rises --- $0.3388$ at $0.5$ and $0.4400$ at $1.0$
against the baseline's $0.6415$ and $0.8940$, a $51\%$ reduction and the
largest effect anywhere in this study. This axis uses $b=15$ with two blocks, a
heavy but still \emph{sparse} corruption of a YaleB image ($20\%$ of pixels);
the explicit $\bS$ absorbs it entirely, whereas every reweighting scheme must
carry some of it in the loss. Where corruption really is sparse and large,
modelling it beats discounting it by a wide margin.

\subsection{Clustering quality}
The two families of metric answer different questions --- RRE asks whether the
\emph{pixels} were recovered, accuracy and NMI ask whether the
\emph{representation} $\bH$ stayed class-discriminative --- so they need not
agree, and it is worth checking whether they do.

On ORL they agree on the ordering but disagree sharply on magnitude. At $b=10$,
$L_1$-NMF improves RRE by $29\%$ ($0.3277 \to 0.2328$) but improves clustering
accuracy by $81\%$ ($0.271 \to 0.491$), with NMI moving from $0.495$ to
$0.667$: the reconstruction metric \emph{understates} what is at stake. RRE
averages over all $1110$ pixels and is dominated by the many every method gets
roughly right, whereas $K$-means on the columns of $\bH$ is sensitive to
whether the occlusion has been allowed to define its own latent directions.
Once basis atoms are spent representing bright squares
(Section~\ref{sec:res-qualitative}), an image's coefficient vector encodes
\emph{where its occlusion landed} as much as \emph{whose face it is}. We report
accuracy under both the majority-vote mapping used in the assignment notebook
and the stricter Hungarian matching; they differ by at most $0.03$ and never
change a ranking.

Extended YaleB produces a genuine \emph{disagreement} between the two metric
families, which is worth reporting rather than smoothing over. At $b=15$ the
sparse-outlier model halves the baseline's RRE ($0.3227$ against $0.6675$) yet
its accuracy, $0.095$ against $0.090$, is indistinguishable. At $b=10$ it is
starker: $L_1$-NMF has by far the best RRE ($0.2682$ against $0.4502$) and the
\emph{worst} accuracy of the five ($0.100$, against $0.108$ for the baseline
and $0.166$ for hypersurface).

On YaleB the dominant source of variance is illumination, not identity, and the
two metrics respond differently. RRE rewards recovering pixel values, which are
dominated by the lighting. Clustering rewards a representation in which images
of the same person land together \emph{despite} the lighting --- a harder,
partly orthogonal requirement --- and the aggressive pixel-level reweighting
that makes $L_1$-NMF a good reconstructor also flattens the identity-carrying
detail $K$-means needs. Hypersurface-cost NMF, with the gentlest reweighting,
is correspondingly the best clusterer among the robust methods at $b=10$. On a
dataset with a strong nuisance factor, a good reconstruction error is not
evidence of a good representation --- and reporting only RRE, which the brief
permits, would have hidden this entirely.

\subsection{Qualitative results}
\label{sec:res-qualitative}
\begin{figure}[tbp]
  \centering
  \resultfig[0.86\linewidth]{orl_reconstruction.png}
  \caption{ORL: the same occluded faces reconstructed by all five algorithms.}
  \label{fig:reconstruction}
\end{figure}
\begin{figure}[tbp]
  \centering
  \resultfig[0.62\linewidth]{orl_basis.png}
  \caption{ORL: leading basis images of $\bW$ learned from occluded data.}
  \label{fig:basis}
\end{figure}
Figures~\ref{fig:reconstruction} and~\ref{fig:basis} show the mechanism behind
the numbers, and they show it unambiguously.

In the learned bases (Figure~\ref{fig:basis}) the Frobenius baseline devotes
several leading atoms to the occlusion itself: two of the six shown are
dominated by a bright patch with visibly \emph{straight, axis-aligned edges},
the signature of a square block rather than of any facial feature.
$L_{2,1}$-NMF's basis is atom for atom the same picture --- the visual
counterpart of its RRE matching the baseline to three decimals.
Hypersurface-cost NMF is intermediate, its patches rounder and softer-edged.
The two pixel-level robust methods show \emph{no} square artefacts at all;
their atoms are grainy but recognisably facial, carrying hairlines, foreheads
and eye regions.

The reconstructions (Figure~\ref{fig:reconstruction}) tell the same story from
the other side. The baseline, hypersurface and $L_{2,1}$ reconstructions all
reproduce a bright blob exactly where the occluding block was, i.e.\ they have
fitted the corruption. $L_1$-NMF and the sparse-outlier model largely
\emph{inpaint} over it, returning a plausible if noisier face. We should be
honest about the limits of this: in the third row the bright region is
attenuated rather than removed, so even the best of the five is suppressing the
occlusion, not detecting and excising it. The grain visible in those two
reconstructions is also the price already quantified in
Section~\ref{sec:res-blocksize} --- the same down-weighting that rejects the
block also discards some genuine high-frequency detail.

\subsection{Convergence and computational cost}
\label{sec:convergence}
\begin{figure}[tbp]
  \centering
  \begin{subfigure}{0.46\linewidth}
    \resultfig[\linewidth]{orl_convergence.png}
    \caption{ORL}
  \end{subfigure}
  \hfill
  \begin{subfigure}{0.46\linewidth}
    \resultfig[\linewidth]{yaleb_convergence.png}
    \caption{Extended YaleB}
  \end{subfigure}
  \caption{Normalised suboptimality against iteration. Each algorithm
           minimises a different objective, so the curves are normalised to
           start at $1$ and are comparable in \emph{rate}, not in level.}
  \label{fig:convergence}
\end{figure}
Three things are worth reporting here, and the third is a limitation rather
than a result.

\paragraph{Monotonicity.} All five objective sequences are non-increasing at
every recorded iteration on both datasets, as the majorise--minimise
derivations guarantee. Since a rising objective would necessarily be a bug, the
check doubles as a correctness test; it is automated in
\texttt{test\_algorithms.py}.

\paragraph{Rate.} Measured as the iteration at which the objective first comes
within $1\%$ of its final value (ORL / YaleB): $L_{2,1}$-NMF $\approx 320/185$,
the baseline $\approx 355/210$, $L_1$-NMF $\approx 390/235$, hypersurface
$\approx 380/240$, the sparse-outlier model $\approx 395/245$. The reweighted
schemes are slower because $\bD$ changes every iteration, so each inner
quadratic is a slightly different problem --- classic IRLS behaviour.

\paragraph{Cost.} Per-iteration cost follows the same ordering and is dominated
by the number of $O(mnk)$ matrix products. On ORL a full $400$-iteration fit
takes $1.8$\,s for the baseline, $3.2$--$3.6$\,s for $L_{2,1}$ and $L_1$, and
$4.1$--$4.5$\,s for hypersurface and the sparse-outlier model; robustness
costs roughly a factor of two in wall-clock time, which is cheap relative to
what it buys in the $b=10$ regime. On Extended YaleB, whose matrix is
$2016 \times 2414$ rather than $1110 \times 400$, a $250$-iteration fit takes
$4.7$--$5.8$\,s for the baseline against $13$--$15$\,s for the sparse-outlier
model, and the full $325$-factorisation grid completes in about an hour on a
modern laptop. The ratio between algorithms is preserved, confirming that the
overhead is the extra matrix products in the reweighting rather than a fixed
cost.

\paragraph{A limitation we should state rather than hide.} With the tolerance
set to $\tau = 10^{-5}$ on the \emph{relative} change of the objective, no run
in the main grid triggered the early stop: every fit used its full iteration
budget, and the \texttt{converged} column of the raw CSVs is $0$ throughout.
The curves in Figure~\ref{fig:convergence} are visibly flat by then --- the
last $10\%$ of iterations move the objective by well under $1\%$ --- so the
comparisons are between fits that are close to a stationary point, but they are
not between fits that have provably stopped moving. We report the budget as
part of the experimental condition for this reason. Multiplicative updates are
in any case guaranteed only to be monotone, and \citet{lin2007convergence}
shows that even convergence to a stationary point requires a modified scheme,
so a tighter tolerance would buy less than it appears to.

\subsection{Hyper-parameter sensitivity}
\label{sec:sensitivity}
\begin{figure}[tbp]
  \centering
  \begin{subfigure}{0.24\linewidth}
    \resultfig[\linewidth]{orl_sensitivity_delta.png}
    \caption{$L_1$ smoothing $\delta$}
  \end{subfigure}
  \hfill
  \begin{subfigure}{0.24\linewidth}
    \resultfig[\linewidth]{orl_sensitivity_scale.png}
    \caption{hypersurface scale $c$}
  \end{subfigure}
  \hfill
  \begin{subfigure}{0.24\linewidth}
    \resultfig[\linewidth]{orl_sensitivity_lam.png}
    \caption{outlier penalty $\lambda$}
  \end{subfigure}
  \hfill
  \begin{subfigure}{0.24\linewidth}
    \resultfig[\linewidth]{orl_sensitivity_rank.png}
    \caption{rank $k$}
  \end{subfigure}
  \caption{ORL: sensitivity of the conclusions to the three free constants.}
  \label{fig:sensitivity}
\end{figure}
Three of our five algorithms carry a free constant, and a comparison whose
conclusions depend on how those constants were chosen would be worthless. We
sweep each of them on ORL at the reference condition $b=10$, one block, all
images corrupted, where the baseline scores $0.3277$.

\paragraph{$L_1$ smoothing $\delta$.} A clear plateau contains the default: RRE
is $0.2404$ at $10^{-4}$ and $0.2310$ at $10^{-3}$ (our default), degrading to
$0.2892$ at $10^{-5}$, where the weights explode on well-fitted pixels, and to
$0.3098$ and $0.3232$ at $10^{-2}$ and $10^{-1}$, where the loss is effectively
quadratic throughout and degenerates towards the baseline. The advantage
survives two full orders of magnitude.

\paragraph{Hypersurface scale $c$.} This sweep documents the trap of
Section~\ref{sec:alg-hypersurface} rather than a plateau. RRE rises
monotonically with $c$: $0.2888$ at $0.01$, $0.3151$ at our default $0.05$,
$0.3202$ at $0.1$, and $0.3290$ at $c=1$ --- the value used in the original
formulation, and within $0.001$ of the Frobenius baseline's $0.3277$. The
degeneration is a measurable fact, not a theoretical worry, so the hypersurface
cost cannot be used off the shelf on data rescaled to $[0,1]$. Our default is
also conservative: $c=0.01$ would have flattered this algorithm, and we keep
$0.05$ because it was fixed before the results were seen.

\paragraph{Outlier penalty $\lambda$.} Both predicted failure modes appear
either side of a working range. RRE is best at $\lambda = 0.02$ ($0.2269$); our
theory-chosen default $1/\sqrt{\max(m,n)} = 0.030$ scores $0.2346$. Below it,
$\lambda = 0.01$ gives $0.2788$ --- $\bS$ starts absorbing genuine facial
detail. Above it the method walks back to the baseline ($0.2984$, $0.3164$,
$0.3286$ at $0.05$, $0.1$, $0.4$), by which point $\bS$ is empty and the model
\emph{is} Frobenius NMF. Picking the default from theory and landing inside the
working range is the outcome we wanted.

\paragraph{Rank $k$.} Across $k \in \{10,20,40,60,80\}$ the ranking is
completely stable --- $L_1$-NMF and the sparse-outlier model near $0.23$ at
every rank, the baseline and $L_{2,1}$ near $0.33$ --- so no conclusion depends
on the choice the brief fixes for us. The interesting detail is at the low end:
at $k=10$ the baseline improves to $0.3216$ and hypersurface to $0.2732$, their
best values in the sweep. A rank too small to represent the occlusion acts as
an implicit regulariser --- a cruder second route to robustness. The
sparse-outlier model trends the other way ($0.2294$ at $k=20$ to $0.2471$ at
$k=80$), since extra capacity in $\bW\bH$ competes with $\bS$ for the same
residual.

\subsection{Discussion and reflection}
\paragraph{Robustness is not a scalar.} The single most useful thing we take
away from this study is that "algorithm A is more robust than algorithm B" is
not a well-formed claim. Each of the five losses discounts residuals at a
particular \emph{granularity} --- per pixel for $L_1$ and hypersurface, per
image for $L_{2,1}$, per entry-with-its-own-variable for the sparse-outlier
model --- and each one helps exactly when the corruption varies at that same
granularity. $L_{2,1}$-NMF is the cleanest demonstration: it is the best method
in the study at a contamination ratio of $0.1$ and numerically identical to the
non-robust baseline at a ratio of $1.0$, without changing a line of its own
code. Had we run only the protocol the brief literally describes --- corrupt
every image, sweep $b$ --- we would have concluded that $L_{2,1}$-NMF simply
does not work, and we would have been wrong for a reason that is invisible from
inside that experiment.

\paragraph{What surprised us.} That the robust losses are clearly \emph{worse}
on clean data ($0.1849$ against $0.1290$ for $L_1$) --- in hindsight just the
efficiency/robustness trade-off, but easy to forget when reading robustness
papers. That the benefit window is narrow at both ends. And that the
hypersurface cost as published is a trap on data scaled to $[0,1]$: with the
paper's $c=1$ every residual stays in the quadratic regime and the method
silently \emph{is} the Frobenius baseline. We found that only because the two
RREs agreed to four decimals --- which is also how we found the $L_{2,1}$
degeneracy. Suspiciously identical numbers were the most productive debugging
signal in the project.

\paragraph{With more compute.} We would run the contamination axis on a finer
grid between $0$ and $0.25$, where the $L_{2,1}$ effect lives and where we
currently have two points; repeat the whole grid at several ranks rather than
only at $k = $ the class count; and replace the fixed iteration budget with a
genuine convergence criterion so that no comparison is confounded by one method
being further from its stationary point than another.

\paragraph{Threats to validity.} (i) A single corruption type: contiguous
saturated blocks. The sparse-outlier model is best-matched to precisely that
noise, so its win is partly a win by construction. (ii) Downsampled greyscale
images, which removes the high-frequency detail a robust loss is most likely to
discard, so our clean-data deficits may be optimistic. (iii) $k$ is fixed to
the class count per the brief, which Section~\ref{sec:sensitivity} shows is not
the best rank for reconstruction. (iv) Multiplicative updates guarantee
monotone descent, not a global optimum, and no run triggered our stopping
tolerance, so all comparisons are between fixed-budget fits. (v) Five
subsamples give coarse error bars (standard deviations of order $0.005$), so
RRE gaps below roughly $0.01$ should not be over-interpreted.

%==============================================================
\section{Conclusion and future work}
\label{sec:conclusion}
% Rubric [3 marks]: "Meaningful conclusions based on results. Meaningful
% future work suggested." Keep it to two short paragraphs and make sure every
% claim is one the experiments actually support.
We implemented five NMF algorithms from scratch in NumPy --- the Frobenius
baseline, $L_1$-norm, hypersurface-cost, $L_{2,1}$-norm and $L_1$-regularised
robust NMF --- derived four of them as instances of one majorise--minimise
template differing only in a weight matrix, and measured them on ORL and
Extended YaleB under occlusion noise varied along three axes, with three
metrics and five $90\%$ subsamples per condition. Three findings stand out.
The explicit sparse-outlier model is the strongest method wherever the
corruption is both heavy and sparse, cutting RRE by $52\%$ against the baseline
at $b=15$ on YaleB and by $51\%$ at full contamination; modelling the
corruption beats discounting it when the noise really is additive and sparse,
which is what a white block is. Second, $L_{2,1}$-NMF is numerically identical
to the non-robust baseline --- to three decimal places, on both datasets ---
whenever every image is corrupted equally, and separates from it only under
uneven contamination, where at a $0.1$ ratio it is the best of the five on ORL
and the best clusterer on YaleB. Robustness is therefore not a scalar property
of an algorithm but a match between the granularity at which the loss discounts
residuals and the granularity at which the corruption varies. Third, RRE and
clustering quality can disagree outright: on YaleB at $b=10$, $L_1$-NMF has the
best reconstruction of the five and the worst clustering, because the
reweighting that rejects the occlusion also flattens the detail that separates
identities under strong illumination variation.

Four directions follow. (i) A \emph{structured} sparse outlier term ---
group-sparse or total-variation-regularised $\bS$ --- would exploit the
contiguity of the occlusion, information the entrywise $L_1$ penalty throws
away; since the sparse-outlier model already wins where it applies, this is the
most promising extension. (ii) A \emph{combined} weighting
$D_{ij} = d_j \cdot w_{ij}$, per-image times per-pixel, should recover the
advantages of $L_{2,1}$ and $L_1$ at once rather than forcing the choice we
observed --- our contamination axis shows the two families win in disjoint
regimes. (iii) Other corruption types would test how much of the sparse-outlier
model's advantage is a win by construction, since our noise matches its
generative assumption exactly. (iv) NNDSVD initialisation
\citep{boutsidis2008svd} and accelerated solvers such as HALS would let every
comparison run to a genuine stationary point rather than to a fixed iteration
budget --- the main methodological limitation of this study.

%==============================================================
\bibliographystyle{plainnat}
\bibliography{references}

%==============================================================
\appendix
\section{Contribution of each group member}
% REQUIRED by the brief ("Indicate the contribution of each group member").
% [H] (not [tbp]) so the table cannot drift out of the appendix into the
% reference list -- an appendix table must stay where it is written.
\begin{table}[H]
  \centering
  \small
  \begin{tabular}{lll}
    \toprule
    Member & Contribution & Approx.\ share \\
    \midrule
    \todo{Name 1} & \todo{e.g.\ algorithm implementation, MM derivations} & \todo{25\%} \\
    \todo{Name 2} & \todo{e.g.\ experiment driver, ORL experiments}       & \todo{25\%} \\
    \todo{Name 3} & \todo{e.g.\ YaleB experiments, figures}               & \todo{25\%} \\
    \todo{Name 4} & \todo{e.g.\ report writing, related work}            & \todo{25\%} \\
    \bottomrule
  \end{tabular}
\end{table}

\section{How to run the code}
% REQUIRED by the brief; omitting it costs 20 marks.
\label{sec:howtorun}
The submitted \texttt{code/} folder is organised as
\begin{verbatim}
code/
  algorithm/        the five NMF implementations + noise, data I/O, metrics
  data/             EMPTY -- copy ORL/ and CroppedYaleB/ in here
  results/          CSVs and figures produced by the scripts
  run_experiment.py run_convergence.py run_sensitivity.py make_figures.py
  run_all.sh        reproduces every number and figure in this report
  test_algorithms.py  self-checks; needs no dataset, runs in under a minute
\end{verbatim}

\paragraph{Environment.}
Python 3.9 or newer. Install the dependencies with
\begin{verbatim}
pip install -r code/requirements.txt
\end{verbatim}
The NMF algorithms themselves import only NumPy and the Python standard
library. SciPy and scikit-learn are used exclusively for evaluation
($K$-means, accuracy, NMI, Hungarian matching); Pillow is used only to decode
the \texttt{.pgm} files and Matplotlib only to draw figures.

\paragraph{Data.}
Copy the two dataset folders so that \texttt{code/data/ORL/} and
\texttt{code/data/CroppedYaleB/} exist.

\paragraph{Reproducing everything.}
\begin{verbatim}
cd code
bash run_all.sh            # ~20 min (ORL) + ~3 h (Extended YaleB)
FAST=1 bash run_all.sh     # reduced grid, a few minutes, for a smoke test
\end{verbatim}

\paragraph{Verifying the implementation.}
\begin{verbatim}
cd code
python test_algorithms.py
\end{verbatim}
This asserts that non-negativity is preserved, that each objective decreases
monotonically (the property the MM derivations of
Section~\ref{sec:methods} guarantee), that every algorithm recovers a
synthetic rank-$k$ non-negative matrix, that the robust losses beat the
Frobenius baseline on corrupted synthetic data, that the noise generator
corrupts exactly the images and pixels it claims, and that all results are
deterministic under a fixed seed. It requires no dataset.

\paragraph{Running a single experiment.}
\begin{verbatim}
cd code
python run_experiment.py  --dataset orl   --data-root data
python run_experiment.py  --dataset yaleb --data-root data --max-iter 200
python run_convergence.py --dataset orl   --data-root data
python run_sensitivity.py --study delta --dataset orl --data-root data
python make_figures.py    --dataset orl   --data-root data
\end{verbatim}
Every script writes its CSVs and PNGs into \texttt{results/} and accepts
\texttt{-{}-help}. All randomness is seeded through \texttt{-{}-seed}
(default \texttt{2026}), so the reported numbers are reproducible.

\section{Additional results}
The main text reports RRE for every condition. The corresponding clustering
accuracy and normalised mutual information tables are produced by
\texttt{code/make\_tables.py} and are available in
\texttt{report/tables/}; the per-run values behind every mean and standard
deviation in this report are in \texttt{code/results/orl\_raw.csv} and
\texttt{code/results/yaleb\_raw.csv}.

\paragraph{Reproducibility.} Every experiment in this report was run twice, on
two different machines with different BLAS implementations (a two-core cloud VM
with OpenBLAS, and an Apple-silicon laptop with Accelerate). The RRE values
agreed to four decimal places on every condition we compared, which is what we
would expect from a deterministic implementation with fixed seeds and is a
useful check that none of the conclusions depends on the numerical
environment.

\end{document}
